\chapter[Sermon 36. Here Is Expounded Who They Be That Are Clothed in White, from Whence Is Salvation, and What Is the True Blessedness.]{Here Is Expounded Who They Be That Are Clothed in White, from Whence Is Salvation, and What Is the True Blessedness.}
\label{sermon:036}
\markright{Sermon 36. Here Is Expounded Who They Be That Are Clothed in ...}

And one of the elders answered, saying to me, “What are these who are arrayed in long white garments, and from where do they come?” And I said to him, “Lord, you know.” And he said to me, “These are those who have come out of great tribulation, and have made their garments clean, and made them white in the blood of the Lamb. Therefore they are in the presence of the throne of God, and serve him day and night in his temple, and he who sits on the throne will dwell among them. They shall hunger no more, neither thirst, neither shall the sun light on them, neither any heat. For the Lamb who is in the midst of the throne shall feed them, and shall lead them to fountains of living water. And God shall wipe away all tears from their eyes.”

\index{Antichrist}John saw the souls of martyrs resting under the altar, clothed in white garments. He also saw an immense multitude from every nation and people, saved from the idolatry of the Gentiles and the superstition of Antichrist, also clothed in white garments. On this occasion, three things shall now be explained to us: what are those who are clothed in white garments? Whence do they have that whiteness, purity, and salvation? Finally, what is the state or felicity of these, or what is true blessedness?

What time John had seen them, he marveled without doubt as to what they were; nevertheless, he seems to have inquired nothing. But of his own accord, one of the twenty-four elders offers himself as an expositor to him, doubtless an excellent teacher, a patriarch and prophet, finally a celestial master, to whom we may justly give credit. Here appears the ignorance of human wit. For just as the eunuch of Ethiopia in chapter 8 of the Acts of the Apostles acknowledges his ignorance, except an interpreter and skilled teacher were given to him, so at this present also blessed John himself, when asked whether he knew those who were clothed in white, confesses his lack of knowledge; yet he ascribes the knowledge to his teacher, by this means requiring a further declaration through a most humble modesty. Finally, here appears the immeasurable goodness of God, which undertakes to teach us who are rude and unworthy. We have many examples of this everywhere in the prophets and in the holy gospel of Christ.

And in the beginning he declares plainly to John, and to all the faithful in the world, what these are who are clothed in white: and he explains fully from whence they came. For with one and the same answer he dispatches both. He says briefly, that the clothed with white in Heaven are the godly people of all times and ages, who at length have escaped out of great tribulation. Tribulation is found to be sundry and diverse. For first, it is tribulation which comes from the plotting and persecution of tyrants. This pertains to martyrs alone. Whereof we have spoken in chapter 6. Which, for as much as in this world they were overwhelmed with unspeakable reproaches for the word of God, they have in another world received white garments. Then there is another tribulation, which arises from the fear of God, and is a care of obtaining salvation. This is sorrowful because of the unrighteousness and corruption of man. It is sorrowful because of the grievous abominations of Antichrist. And these also, albeit they be not made martyrs, yet are they in another life clothed with white. Finally, they have tribulation and are molested in the flesh, so many as mortify their flesh with its lusts. And because they mourn here in the world, they shall receive comfort and consolation in the world to come.

\index{Arethas of Caesarea}\index{Primasius}\index{Repentance}\index{Hebrews, Epistle to the}Again, lest any man should ascribe life and salvation to martyrdom as to our work, and to repentance as to our desert, the Lord moreover declares expressly from whence that life and salvation proceeds, and how that whiteness and purity comes to us. And they have washed their garments, he says, and made them white by the blood of the Lamb. And here is found a diverse reading. Some suggest that they have enlarged their garments, so that it might appear he alluded to the families of great princes, who use, for the setting forth of their renown, to put on most wide and most sumptuous garments. But in my opinion, the Complutensian copy and Arethas seem to read more rightly and more simply, namely, they have washed, as also the old interpreter has translated it. For by this exposition follows, and have made white. Primasius reads, and have made their garments white in the blood of the Lamb. And hereby is signified that the salvation and purification of the faithful is of the blood of Christ, and of no other thing. Where verily blood sprinkled whitens not, but pollutes. Therefore we must understand these things spiritually, to wit that the very, natural and human blood of Christ, shed upon the cross, being sprinkled upon us spiritually (as Paul to the Hebrews 10 explains), and received with faith, although it touches us not naturally and corporally, purges us from all sins. And therefore we read in another place that Christ purges us with his blood. For because sanctification is the only work of God, therefore where the saints are said now to have washed and whited their garments by the blood of the Lamb, it signifies that they have received by faith the purification prepared by the blood. And this doctrine is catholic and of the right faith, which has so many and so great testimonies in the holy Scriptures. Finally, we perceive how those who are saved from the kingdom of Antichrist are saved by the merit of Christ alone, and by no other thing, as I have also shown you before. Moreover, it follows immediately: therefore they are in the sight of God’s seat. For what cause, I pray thee? Because they have washed and whited their garments in the blood of the Lamb: therefore for the merit of Christ have they entered into heaven, and there are surrounded with eternal light.

Finally, the elder proclaims at length the condition of the saints and the true blessedness of the faithful. These things are certain glimpses, provided here for consolation. Or else, such things as the eye has not seen, or the ear heard, that God has prepared for those who love him. He recounts many things, from which we are to understand the excellence of eternal salvation and what good things we attain in it.

First, the saints stand before the throne of God. In the throne is the majesty of God to be worshipped forever, and the blessed Trinity. And the saints stand before the seat, not as those who tarry before the gates. For as God’s most complete friends, they are always in God’s sight and have the enjoyment of his deity. Concerning this, the Lord speaks in the Gospel: “Pray that you may escape these things and stand before the Son of Man.” And also David: “The fulfillment of joys is in your sight, and pleasure at your right hand forever.” And there is added another thing, which may explain that standing: they serve God in his temple both day and night. That service has pleasure and no pain. And they serve God in the temple, as God is wont to be served in the temple. For they keep holy days, are glad, rejoice, are merry, praise, and so they offer up sacrifices and are refreshed with heavenly repast. And this joy shall be everlasting and perpetual, which is signified by day and night. Otherwise, in the blessedness everlasting there is no might at all, nor any changeable course of time. Hereunto is added that he who sits in the seat, that is, the divine majesty, [?will] dwell in them: that is to say, God will be all in all, or he will lean over them, and as it were a tent or tabernacle, will overshadow them, defend and keep them, and give himself whole to be enjoyed by them, as most familiar and friendly to them. Moreover, they shall hunger no more, neither shall they thirst. For all infirmity and misery is taken away from the blessed souls and glorified bodies. They are filled with all good things without any loathsomeness, with a most joyous fulfillment. Now the sun falls not upon them, nor the heat: which phrase of speech betokens that they are put to no labor nor pain, but are delivered at once from all displeasure and all pain, and to be at most pleasant rest.

\index{Ezekiel (Book of)}Again, this is established for so great a blessedness, Christ the Lamb—that is to say, Christ the mediator and redeemer—in the midst of the throne, that is, very God. For he, as both Ezekiel 34:23 and the Lord himself in John 10:11 testify, will feed them like a shepherd, and as a leader of life will lead them to the fountains of living water: that is to say, he will quicken them forever and preserve all his in that blessedness. He uses in this treatise words of the prophets most commonly used, so that, climbing to higher things, we might in some measure esteem heavenly gifts. Hereunto he joins yet a notable blessing: and the Lord will wipe away all tears from their eyes. Which words he has borrowed from Isaiah. For saints in this world, troubled by various evils, have shed most plentiful tears; but in the world to come, the Lord will comfort them, gladdening them with everlasting joy, never giving them any occasion of grief. And therefore he said in the Gospel, “Verily I say unto you, you shall weep and lament, but again the world shall rejoice; you shall mourn, but your mourning shall be turned into joy. And your heart shall rejoice, and your joy no one shall take from you.” We shall hear similar things in Revelation 21 as well.

Hereby they perceive how shamefully they sin, those who constantly say, “If I should endure this present life for the sake of religion, who will tell me what that other life to come will be like?” Perhaps if I neglect this, I will gain nothing in another world. For here we have a most manifest testimony that, as most assured salvation is prepared by God in heaven for the faithful, so it is also most ample and great; in so much that the Apostle in another place says that the afflictions of this present time are not equal to the glory which shall be revealed to us. The Lord grant us that we may acknowledge these things.
